\documentclass[12pt,a4paper]{article}

\usepackage{amsmath, amsthm, amssymb}
\usepackage{mathtools}
\usepackage{hyperref}
\hypersetup{
  pdftitle={Semiclassical Consistency of the Weil Representation on Heis3(Z/qZ): Quantitative Sinc Embedding, Generator Convergence, and Aliasing Control},
  pdfauthor={Jérôme Beau},
  pdfsubject={Cosmochrony emergent-geometry programme, paper H2, version 1.1},
}
\newcommand{\doi}[1]{\href{https://doi.org/#1}{doi:#1}}

\usepackage[T1]{fontenc}
\usepackage[utf8]{inputenc}

\newtheorem{theorem}{Theorem}[section]
\newtheorem{lemma}[theorem]{Lemma}
\newtheorem{corollary}[theorem]{Corollary}
\newtheorem{definition}[theorem]{Definition}
\newtheorem{remark}[theorem]{Remark}
\newtheorem{hypothesis}[theorem]{Hypothesis}

\newcommand{\Heis}{\mathrm{Heis}_3}
\newcommand{\Zq}{\mathbb{Z}/q\mathbb{Z}}
\newcommand{\ltwo}{L^2(\mathbb{R})}
\newcommand{\ltwod}{\ell^2(\Zq)}
\newcommand{\Schwartz}{\mathcal{S}(\mathbb{R})}
\newcommand{\iotaq}{\iota_q}
\newcommand{\iotastar}{\iota_q^*}
\newcommand{\Xhat}{\hat{X}_q}
\newcommand{\Phat}{\hat{P}_q}
\newcommand{\norm}[1]{\left\|#1\right\|}
\newcommand{\abs}[1]{\left|#1\right|}

\title{Semiclassical Consistency of the Weil Representation on
  $\Heis(\mathbb{Z}/q\mathbb{Z})$:\\
  Quantitative Sinc Embedding, Generator Convergence,\\
  and Aliasing Control}
\author{J\'er\^ome Beau\\
\small Independent Researcher\\
\small \texttt{jerome.beau@cosmochrony.org}}
\date{2026}

\begin{document}

  \maketitle

  \begin{abstract}
    We prove that the rescaled Weil generators $\Xhat$ and $\Phat$ of the discrete Heisenberg group
    $\Heis(\Zq)$ converge strongly to the position operator $x$ and the momentum operator $-i\partial_x$
    as $q \to \infty$.
    This is the statement formulated as hypothesis~[H2] in version~2.0 of the Cosmochrony paper
    $Q5a$; the current version of $Q5a$ does not use that framework.
    The proof rests on two structural lemmas.
    First, a discrete Sobolev estimate controls the error between the finite-difference operator
    and the continuous derivative at the $\ell^2$ level.
    Second, a quantitative Poisson aliasing lemma controls the transition from
    $\ell^2(\Zq)$ to $\ltwo$ with an explicit rate $O(q^{-1})$ for data in $\Schwartz$.
    As a consequence we obtain a quasi-isometry lemma for the sinc embedding $\iotaq$
    with the explicit rate
    \[
      \abs{\norm{\iotaq f}_{L^2}^2 - \norm{f}_{\ell^2}^2} \leq \frac{C}{q}\,\mathcal{N}(\psi),
    \]
    which is a reusable structural estimate; an identification of an effective metric constant
    $A$ in terms of the Born--Infeld saturation constant $c_{\mathrm{BI}}$ would use it together
    with inputs that are not established (W1).
    In the framework of $Q5a$ version~2.0, [H2] is the generator-convergence hypothesis of its
    large-$q$ limit theorem.
    The other hypotheses of that theorem are not discharged: the common weight hypothesis [H-w]
    is open (W1), and the Q5a-O2 Fourier analysis rules out
    [H-E1] at the $q^{-1}$ scale and supersedes [C].
    For uniformly bounded weights, the current $Q5a$ proves that the published admissibility form
    has zero-form Mosco limit on the canonical filtration, so the present result supplies no
    spatial limit operator.
  \end{abstract}

  \medskip
  \noindent\textbf{Keywords.}
  Weil representation; Heisenberg group; semiclassical limit; sinc embedding;
  discrete Sobolev estimate; aliasing control; Schwartz convergence; Cosmochrony; $Q5a$.

  \newpage
  \tableofcontents

  \newpage
  \section{Introduction}
    \label{sec:intro}

    \subsection{Context and motivation}

      The Cosmochrony programme derives physical observables from the spectral structure
      of admissible non-injective projections on discrete Heisenberg groups
      $\Heis(\Zq)$~\cite{Beau2026m,Beau2026HeisenbergStructure}.
      Version~2.0 of the $Q5a$ paper~\cite{Beau2026q5a} (Zenodo record 19897055) formulated a
      large-$q$ limit theorem for this structure (Theorem~7.1 there), identifying the limiting
      Hilbert space as $\ltwo$ and the limiting operators as the standard generators of the
      Schr\"odinger (metaplectic) representation of $\Heis(\mathbb{R})$, under the hypotheses
      [H-$\Pi$1], [H-$\Pi$2], [H2], [H-w] and [H-E1], together with a spectral tightness
      statement [C] for its form-convergence part.
      That framework is not the one of the current $Q5a$, which proves that the canonical
      filtration is a Fourier window on which, for uniformly bounded weights, the published
      admissibility form has zero-form Mosco limit (\cite{Beau2026q5a}, Theorems~3.1 and~5.1).
      Of the hypotheses of the version-2.0 theorem, the common weight hypothesis [H-w] is open,
      W1~\cite{Beau2026w1} establishing no proof of it, and the Q5a-O2 Fourier
      analysis~\cite{Beau2026q5ao2} rules out [H-E1] at the $q^{-1}$ scale and supersedes [C].
      The present paper addresses one hypothesis of that framework, the generator-convergence
      hypothesis~[H2] (Hypothesis~4.3 of $Q5a$ version~2.0):

      \begin{hypothesis}[\textup{[H2]}, $Q5a$ version~2.0~\cite{Beau2026q5a}]
        \label{hyp:H2}
        For every $\psi \in \Schwartz$ and every $\varepsilon > 0$, there exists $q_0$ such that
        for all primes $q \geq q_0$,
        \begin{align*}
          \norm{\iotaq\bigl(\Xhat\,\iotastar\psi\bigr) - x\psi}_{L^2(\mathbb{R})} &< \varepsilon, \\
          \norm{\iotaq\bigl(\Phat\,\iotastar\psi\bigr) - (-i\partial_x\psi)}_{L^2(\mathbb{R})} &< \varepsilon.
        \end{align*}
        The convergence holds uniformly for $\psi$ in bounded subsets of $\Schwartz$.
      \end{hypothesis}

      The present paper proves the statement of Hypothesis~\ref{hyp:H2}.
      It discharges no hypothesis of the current $Q5a$.

    \subsection{Structure of the argument}

      The two operators $\Xhat$ and $\Phat$ have asymmetric proof structures.
      For the position generator $\Xhat$, the argument is quasi-pointwise:
      the character $e^{2\pi i k/q}$ is approximated by $1 + 2\pi ik/q$ via a Taylor expansion,
      and the error is controlled by the Schwartz decay of $\iotastar\psi$.
      For the momentum generator $\Phat$, the argument is genuinely non-local:
      a discrete Sobolev estimate (Lemma~\ref{lem:sobolev}) is required
      to control the finite-difference error at the $\ell^2$ level before the aliasing
      lemma (Lemma~\ref{lem:aliasing}) can be applied to pass to $\ltwo$.
      The logical chain for $\Phat$ is therefore:
      \begin{center}
        discrete Sobolev estimate $\longrightarrow$ aliasing lemma $\longrightarrow$ $\ltwo$ convergence.
      \end{center}
      This asymmetry is reflected explicitly in the proof structure of
      Sections~\ref{sec:position} and~\ref{sec:momentum}.

    \subsection{Relation to previous work}

      Hannay and Berry~\cite{HannayBerry1980} study the semiclassical limit of quantised maps
      on the torus in the dual regime $\hbar = (2\pi q)^{-1} \to 0$; their limit shares the same
      discrete-to-continuous mechanism.
      Degli Esposti and Graffi~\cite{DegliEspostiGraffi1993} provide structural support for~[H2]
      via an analysis of the semiclassical limit of maps quantised on $\mathbb{Z}/q\mathbb{Z}$,
      but their results do not cover $\Heis(\Zq)$ in the parameterisation of~\cite{Beau2026m}.
      The present paper fills this gap.

    \subsection{Reusability of the quasi-isometry lemma}

      The quasi-isometry bound of Corollary~\ref{cor:quasi-iso} (Section~\ref{sec:aliasing}) is not
      merely a technical step in the proof of~[H2].
      It quantifies the defect between $\norm{\iotaq f}_{L^2}^2$ and $\norm{f}_{\ell^2}^2$
      with a constant that depends explicitly on Sobolev norms of $\psi$ but is independent of $q$.
      In $Q5a$ version~2.0, open problem $Q5a$-O5 asked to express the constant $A > 0$ of the limit
      form $\mathcal{E}(f,g) = A\int_{\mathbb{R}} \partial_x f\,\partial_x g\,dx$ in terms of the
      Born--Infeld saturation constant and the BFS combinatorial data of $G_q$.
      Version~1 of W1 proposed the explicit expression
      \[
        A = \sum_{n \geq 1} \sigma^*(n) \approx \frac{c_{\mathrm{BI}}^2}{C_{\mathrm{Heis}}} \,
        \zeta\!\left(\tfrac{\delta^*}{2}\right);
      \]
      W1~\cite{Beau2026w1} retains the approximation for $A$ only as that proposal, under a
      supplied global power-law profile, saturation and normalisation, and it is not a
      consequence of O25, $Q5a$ or U1 (W1, Section~4.2).
      An identification of $A$ would use a quasi-isometry bound of the present kind as one
      ingredient; this paper supplies the bound, not the identification.
      We therefore state Corollary~\ref{cor:quasi-iso} as a numbered result.

    \subsection{Organisation}

      Section~\ref{sec:embedding} defines the sinc embedding and its adjoint,
      fixes normalisations, and records the uniform operator bound on the Weil generators
      (Lemma~\ref{lem:uniform-bound}).
      Section~\ref{sec:sobolev} proves the discrete Sobolev estimate (Lemma~\ref{lem:sobolev}).
      Section~\ref{sec:aliasing} proves the quantitative aliasing lemma (Lemma~\ref{lem:aliasing})
      and derives the quasi-isometry corollary (Corollary~\ref{cor:quasi-iso}).
      Section~\ref{sec:position} proves convergence of $\Xhat$.
      Section~\ref{sec:momentum} proves convergence of $\Phat$.
      Section~\ref{sec:schwartz} assembles the strong convergence theorem on $\Schwartz$.
      Section~\ref{sec:extension} extends the result to $\ltwo$ by density.

      \newpage
  \section{Sinc embedding and normalisations}
    \label{sec:embedding}

    \subsection{The sinc embedding}

      \begin{definition}[Sinc embedding]
        \label{def:sinc}
        Let $q$ be a positive integer.
        The \emph{sinc embedding} $\iotaq : \ltwod \to \ltwo$ is defined by
        \[
          (\iotaq f)(x) := q^{1/2} \sum_{k \in \Zq} f(k)\,\mathrm{sinc}\!\left(qx - k\right),
          \qquad \mathrm{sinc}(t) := \frac{\sin(\pi t)}{\pi t}.
        \]
      \end{definition}

      \begin{definition}[Sampling operator]
        \label{def:sampling}
        The \emph{sampling operator} $\iotastar : \ltwo \to \ltwod$ is the $L^2$-adjoint of $\iotaq$,
        given by
        \[
          (\iotastar \psi)(k) := q^{1/2}\,\psi\!\left(\tfrac{k}{q}\right), \qquad k \in \Zq.
        \]
      \end{definition}

      \begin{remark}
        \label{rem:normalisation}
        The factor $q^{1/2}$ in Definition~\ref{def:sinc} ensures that
        $\iotaq \iotastar$ is approximately the identity on the low-frequency sector of $\ltwo$,
        consistently with the Nyquist--Shannon sampling theorem.
        The factor $q^{1/2}$ in Definition~\ref{def:sampling} is the adjoint normalisation.
        $Q5a$ version~2.0 (Remark~3.2 there) writes the sinc interpolation without the factor
        $q^{1/2}$ and describes it as an isometry onto the band-limited subspace; that isometry
        requires the factor $q^{1/2}$ of Definition~\ref{def:sinc}, and hypothesis~[H2] is read
        here with this normalisation, under which it is proved.
      \end{remark}

    \subsection{Weil generators and rescaled generators}

      The Weil generators of $\Heis(\Zq)$ acting on $\ltwod$ are
      \begin{align*}
        (\rho_q(X)f)(k) &:= e^{2\pi ik/q} f(k), \\
        (\rho_q(Y)f)(k) &:= f(k+1 \bmod q).
      \end{align*}
      The \emph{rescaled generators} are defined as in $Q5a$ version~2.0~\cite{Beau2026q5a}
      (Definition~4.1 there):
      \[
        \Xhat := \frac{q}{2\pi i}\bigl(\rho_q(X) - I\bigr), \qquad
        \Phat := \frac{q}{2\pi i}\bigl(\rho_q(Y) - I\bigr).
      \]

    \subsection{Uniform operator bound}

      The following lemma is used in Section~\ref{sec:extension} to extend
      the $\Schwartz$ convergence result to $\ltwo$ by density.

      \begin{lemma}[Uniform operator bound]
        \label{lem:uniform-bound}
        For all primes $q$ and both generators,
        \[
          \norm{\rho_q(X) - I}_{\ltwod \to \ltwod} \leq 2, \qquad
          \norm{\rho_q(Y) - I}_{\ltwod \to \ltwod} \leq 2.
        \]
        Consequently,
        \[
          \norm{\Xhat}_{\ltwod \to \ltwod} \leq \frac{q}{\pi}, \qquad
          \norm{\Phat}_{\ltwod \to \ltwod} \leq \frac{q}{\pi}.
        \]
        After transport by $\iotaq$, the operators $\iotaq \Xhat \iotastar$ and
        $\iotaq \Phat \iotastar$ are bounded on $\ltwo$ uniformly in $q$ on the range of $\iotastar$.
        Although the $\ltwod$-operator norms $\norm{\Xhat}$, $\norm{\Phat}$ grow as $O(q)$,
        the action on sampled Schwartz data $\iotastar\psi$ is effectively controlled via the quasi-isometry
        (Corollary~\ref{cor:quasi-iso}): the transported operators $\iotaq\Xhat\iotastar$ and
        $\iotaq\Phat\iotastar$ remain bounded uniformly in $q$ on bounded subsets of $\Schwartz$,
        since the Schwartz decay of $\psi$ absorbs the $O(q)$ growth.
      \end{lemma}

      \begin{proof}
        [Proof sketch]
Both $\rho_q(X)$ and $\rho_q(Y)$ are unitary on $\ltwod$, so $\norm{\rho_q(s) - I} \leq 2$
for $s \in \{X, Y\}$.
The bound on the rescaled generators follows from the definition and the triangle inequality.
      \end{proof}

      \newpage
  \section{Discrete Sobolev estimate}
    \label{sec:sobolev}

    Let $T_1$ denote the shift operator on $\ltwod$: $(T_1 g)(k) := g(k+1 \bmod q)$,
    and let $T_{-1}$ denote its inverse.
    Define the \emph{forward discrete derivative} $D_q g := (q/2\pi i)(T_1 g - g)$,
    which coincides with $\Phat g$ by definition, and the \emph{centered discrete derivative}
    \[
      g'_{\mathrm{disc}} := \frac{q}{4\pi i}(T_1 - T_{-1})g,
    \]
    which is the standard centered finite-difference approximation of the derivative at step size $1/q$.

    \begin{lemma}[Discrete algebraic identity]
      \label{lem:sobolev}
      For every $g : \Zq \to \mathbb{C}$, the following identity holds exactly:
      \[
        D_q g - g'_{\mathrm{disc}}
        = \frac{q}{4\pi i}(T_1 - 2I + T_{-1})g.
      \]
      In particular,
      \[
        \norm{D_q g - g'_{\mathrm{disc}}}_{\ell^2}
        = \frac{q}{4\pi}\,\norm{(T_1 - 2I + T_{-1})g}_{\ell^2}.
      \]
    \end{lemma}

    \begin{proof}
      [Proof sketch]
Direct computation: $(T_1 - I) = \frac{1}{2}(T_1 - T_{-1}) + \frac{1}{2}(T_1 - 2I + T_{-1})$.
Multiplying both sides by $q/(2\pi i)$ gives the identity.
    \end{proof}

    \begin{lemma}[Passage to the continuous second derivative]
      \label{lem:sobolev-continuous}
      For $g = \iotastar\psi$ with $\psi \in \Schwartz$,
      \[
        \norm{(T_1 - 2I + T_{-1})\iotastar\psi}_{\ell^2} \leq C'\,\norm{\partial_x^2 \psi}_{L^2(\mathbb{R})}.
      \]
    \end{lemma}

    \begin{proof}
      [Proof sketch]
This follows from the aliasing lemma (Lemma~\ref{lem:aliasing} below) applied to
the sampled second finite difference.
    \end{proof}

    \begin{remark}
      Lemma~\ref{lem:sobolev} is purely $\ell^2$-intrinsic and does not reference any
      continuous norm.
      The passage to $\norm{\partial_x^2\psi}_{L^2}$ in Lemma~\ref{lem:sobolev-continuous}
      requires the aliasing control of Section~\ref{sec:aliasing}.
      This creates a logical dependency: Section~\ref{sec:sobolev} is independent of
      Section~\ref{sec:aliasing} at the lemma level, but the application to
      $g = \iotastar\psi$ (used in Section~\ref{sec:momentum}) requires both.
    \end{remark}

    \newpage
  \section{Uniform aliasing control}
    \label{sec:aliasing}

    \subsection{Poisson aliasing lemma}

      \begin{lemma}[Quantitative aliasing bound]
        \label{lem:aliasing}
        For $\psi \in \Schwartz$ and any $M \geq 2$, there exists $C_M > 0$ depending only on
        $\psi$ and $M$ such that
        \[
          \sum_{k \neq 0} \abs{\hat{\psi}(\xi + kq)} \leq \frac{C_M}{q^M}
        \]
        uniformly in $\xi \in \mathbb{R}$.
        Consequently, for the integral aliasing error,
        \[
          \abs{\int_{\mathbb{R}} \hat{\psi}(\xi + kq)\,\overline{\hat{\psi}(\xi)}\,d\xi}
          \leq \frac{C_M'}{q^M}
        \]
        for each $k \neq 0$.
      \end{lemma}

      \begin{proof}
        [Proof sketch]
Since $\psi \in \Schwartz$, its Fourier transform satisfies
$\abs{\hat{\psi}(\xi)} \leq C_M\,(1 + \abs{\xi})^{-M}$ for all $M$, uniformly over all $\xi \in \mathbb{R}$
(this polynomial decay is the key property that makes the bound uniform in $\xi$).
For $k \neq 0$ and $q \geq 1$, $\abs{\xi + kq} \geq \abs{k}q - \abs{\xi}$,
so for $\abs{\xi} \leq q/2$ one has $\abs{\xi + kq} \geq (\abs{k} - \tfrac{1}{2})q$.
Summing over $k \neq 0$ and bounding $\sum_{k \neq 0} \abs{k}^{-M} < \infty$ for $M > 1$
yields the first bound.
The estimate extends to all $\xi \in \mathbb{R}$ by periodicity of the aliasing decomposition.
The integral estimate follows by Cauchy--Schwarz.
      \end{proof}

    \subsection{Quasi-isometry with explicit rate}

      \begin{corollary}[Quasi-isometry with rate $O(q^{-1})$]
        \label{cor:quasi-iso}
        For $\psi \in \Schwartz$ and $f = \iotastar\psi$,
        \[
          \abs{\norm{\iotaq f}_{L^2}^2 - \norm{f}_{\ell^2}^2}
          \leq \frac{C}{q}\,\mathcal{N}_2(\psi),
        \]
        where $\mathcal{N}_2(\psi) := \norm{x\psi}_{L^2}^2 + \norm{\partial_x\psi}_{L^2}^2$
        and $C > 0$ is independent of $q$.
      \end{corollary}

      \begin{proof}
        [Proof sketch]
The Parseval identity for the sinc embedding gives
\[
  \norm{\iotaq f}_{L^2}^2 - \norm{f}_{\ell^2}^2
  = \sum_{k \neq 0} \int_{\mathbb{R}} \hat{\psi}(\xi + kq)\,\overline{\hat{\psi}(\xi)}\,d\xi.
\]
Applying Lemma~\ref{lem:aliasing} with $M = 2$ and bounding the sum over $k \neq 0$
by a geometric series in $q$ yields the $O(q^{-1})$ rate.
The explicit dependence on $\mathcal{N}_2(\psi)$ comes from the $M = 2$ decay estimate
applied to the $k = \pm 1$ dominant terms.
      \end{proof}

      \begin{remark}[Reusability for an identification of $A$]
        \label{rem:reusability}
        The rate $O(q^{-1})$ in Corollary~\ref{cor:quasi-iso} is summable with respect to
        weights $\sigma^*(n) \propto n^{-\delta^*/2}$ provided $\delta^*/2 > 1$.
        Such a global power-law profile is a supplied input, not an established
        one~\cite{Beau2026w1}.
        Under it, the bound is one quantitative ingredient that an identification of $A$ in terms
        of $c_{\mathrm{BI}}$ would use; it does not supply that identification.
      \end{remark}

      \newpage
  \section{Convergence of the position generator}
    \label{sec:position}

    \begin{theorem}[Convergence of $\Xhat$]
      \label{thm:position}
      For every $\psi \in \Schwartz$,
      \[
        \norm{\iotaq\bigl(\Xhat\,\iotastar\psi\bigr) - x\psi}_{L^2(\mathbb{R})}
        = O\!\left(q^{-1}\right),
      \]
      uniformly on bounded subsets of $\Schwartz$.
    \end{theorem}

    \begin{proof}
      [Proof sketch]
Since $(\Xhat \iotastar\psi)(k) = \frac{q}{2\pi i}(e^{2\pi ik/q} - 1)(\iotastar\psi)(k)$,
the Taylor expansion $e^{2\pi ik/q} = 1 + \frac{2\pi ik}{q} + O((k/q)^2)$ gives
\[
  (\Xhat \iotastar\psi)(k) = k\cdot(\iotastar\psi)(k) + E_q(k),
  \qquad \abs{E_q(k)} \leq \frac{C k^2}{q} \abs{(\iotastar\psi)(k)}.
\]
Step~1: The main term.
One shows $\norm{\iotaq(k \cdot \iotastar\psi) - x\psi}_{L^2} \leq C'/q$ using
the aliasing lemma (Lemma~\ref{lem:aliasing}) to compare the sinc reconstruction
of $k \cdot (\iotastar\psi)(k)$ with multiplication by $x$.

Step~2: The remainder.
$\norm{\iotaq E_q}_{L^2} \leq (C/q)\norm{\iotaq(k^2 \iotastar\psi)}_{L^2}
\lesssim (C/q)\norm{x^2 \psi}_{L^2} = O(q^{-1})$
by the quasi-isometry bound (Corollary~\ref{cor:quasi-iso}).
    \end{proof}

    \newpage
  \section{Convergence of the momentum generator}
    \label{sec:momentum}

    \begin{theorem}[Convergence of $\Phat$]
      \label{thm:momentum}
      For every $\psi \in \Schwartz$,
      \[
        \norm{\iotaq\bigl(\Phat\,\iotastar\psi\bigr) - (-i\partial_x\psi)}_{L^2(\mathbb{R})}
        = O\!\left(q^{-1}\right),
      \]
      uniformly on bounded subsets of $\Schwartz$.
    \end{theorem}

    \begin{proof}
      [Proof sketch]
The argument proceeds in three steps.

Step~1: Discrete algebraic decomposition ($\ell^2$ level).
By Lemma~\ref{lem:sobolev}, $\Phat\,\iotastar\psi = g'_{\mathrm{disc}} + \frac{q}{4\pi i}(T_1 - 2I + T_{-1})\iotastar\psi$,
where $g'_{\mathrm{disc}} = \frac{q}{4\pi i}(T_1 - T_{-1})\iotastar\psi$ is the centered discrete derivative.
By Lemma~\ref{lem:sobolev-continuous} (which uses the aliasing lemma), the error term satisfies
\[
  \left\|\frac{q}{4\pi i}(T_1 - 2I + T_{-1})\iotastar\psi\right\|_{\ell^2}
  \leq \frac{C}{q}\,\norm{\partial_x^2\psi}_{L^2}.
\]

Step~2: Approximation of the centered difference by the sampled continuous derivative.
The centered difference $g'_{\mathrm{disc}}(k)$ approximates
the \emph{sampled continuous derivative}
$(\partial_x\psi)|_{\mathbb{Z}/q\mathbb{Z}}$, i.e.\ the sequence $k \mapsto q^{1/2}(-i\partial_x\psi)(k/q)$
obtained via the adjoint embedding $\iotastar$ applied to $-i\partial_x\psi$,
up to an aliasing error controlled by Lemma~\ref{lem:aliasing}:
\[
  \norm{g'_{\mathrm{disc}} - (-i\partial_x\psi)|_{\mathbb{Z}/q\mathbb{Z}}}_{\ell^2}
  \leq \frac{C'}{q}\,\norm{\partial_x^3\psi}_{L^2}.
\]
The third-order derivative in the bound arises from the consistency error of the
centered difference scheme: the leading term of the Taylor remainder for the centered
approximation is of order $(\Delta x)^2 \partial_x^3\psi \sim q^{-2}\partial_x^3\psi$,
which after the $\ell^2$-norm over $q$ points and the factor $q$ from the rescaling
gives $O(q^{-1})\|\partial_x^3\psi\|_{L^2}$.
Combining Steps~1 and~2, one obtains
$\norm{\Phat\,\iotastar\psi - (-i\partial_x\psi)|_{\mathbb{Z}/q\mathbb{Z}}}_{\ell^2} \leq (C + C')/q
\cdot \norm{\psi}_{W^{3,2}}$.

Step~3: Passage from $\ell^2$ to $L^2$ via aliasing.
The aliasing lemma (Lemma~\ref{lem:aliasing}) and the quasi-isometry
(Corollary~\ref{cor:quasi-iso}) give
\[
  \norm{\iotaq(\Phat\,\iotastar\psi) - (-i\partial_x\psi)}_{L^2}
  \leq \norm{\iotaq(\Phat\,\iotastar\psi - (\partial_x\psi)|_{\Zq})}_{L^2}
  + \norm{\iotaq((\partial_x\psi)|_{\Zq}) - \partial_x\psi}_{L^2},
\]
and both terms are bounded by $O(q^{-1})\norm{\partial_x^2\psi}_{L^2}$.
    \end{proof}

    \begin{remark}
      The asymmetry between Theorems~\ref{thm:position} and~\ref{thm:momentum}
      is structural, not a proof artefact.
      The character $e^{2\pi ik/q}$ is a scalar multiplier; its Taylor expansion is local.
      The shift $T_1$ is a non-local operator; its approximation by a derivative requires
      a genuine discrete Sobolev estimate.
    \end{remark}

    \newpage
  \section{Strong convergence on Schwartz data}
    \label{sec:schwartz}

    \begin{theorem}[Strong convergence on $\Schwartz$]
      \label{thm:schwartz}
      For every $\psi \in \Schwartz$,
      \[
        \norm{\iotaq(\Xhat\,\iotastar\psi) - x\psi}_{L^2(\mathbb{R})} \xrightarrow{q \to \infty} 0,
        \qquad
        \norm{\iotaq(\Phat\,\iotastar\psi) + i\partial_x\psi}_{L^2(\mathbb{R})} \xrightarrow{q \to \infty} 0.
      \]
      The convergence is uniform on bounded subsets of $\Schwartz$.
      This is the statement of hypothesis~\textup{[H2]} of $Q5a$ version~2.0~\cite{Beau2026q5a}.
    \end{theorem}

    \begin{proof}
      Immediate from Theorems~\ref{thm:position} and~\ref{thm:momentum}.
    \end{proof}

    \begin{remark}[Relation to the large-$q$ theorem of $Q5a$ version~2.0]
      \label{rem:Q5a-consequence}
      Theorem~7.1 of $Q5a$ version~2.0~\cite{Beau2026q5a} assumes [H-$\Pi$1], [H-$\Pi$2], [H2],
      [H-w] and [H-E1], and its form-convergence part additionally assumes a spectral tightness
      statement [C].
      Theorem~\ref{thm:schwartz} proves the statement of [H2], which is the generator statement of
      that theorem's representational part (T2).
      It discharges no other hypothesis: [H-w] is open, W1~\cite{Beau2026w1} establishing no
      proof of it; the Q5a-O2 Fourier analysis~\cite{Beau2026q5ao2} rules out [H-E1] at the
      $q^{-1}$ scale and supersedes [C]; and the status of [H-$\Pi$1] and [H-$\Pi$2] is not examined
      here.
      The form-convergence part (T3) is not available: for uniformly bounded weights, the current
      $Q5a$ proves that the published admissibility forms have zero-form Mosco limit on the
      canonical Fourier filtration (\cite{Beau2026q5a}, Theorem~5.1), and under its depth and weight hypotheses no common
      scalar normalisation yields a toric differential limit (\cite{Beau2026q5a}, Theorem~6.1);
      a non-trivial limit form $\mathcal{E}(f,g) = A\int_{\mathbb{R}} \partial_x f\,
      \partial_x g\,dx$ is therefore not obtained by that route.
      The framework of $Q5a$ version~2.0 is not that of the current $Q5a$, and no hypothesis of the
      current $Q5a$ is discharged here.
    \end{remark}

    \newpage
  \section{Extension to $L^2(\mathbb{R})$}
    \label{sec:extension}

    \begin{theorem}[Extension to $L^2$]
      \label{thm:extension}
      The strong convergence of Theorem~\ref{thm:schwartz} extends to all $\psi \in \ltwo$:
      for every $\psi \in \ltwo$ and $\varepsilon > 0$, there exists $q_0$ such that
      for all $q \geq q_0$,
      \[
        \norm{\iotaq(\Xhat\,\iotastar\psi) - x\psi}_{L^2} < \varepsilon,
        \qquad
        \norm{\iotaq(\Phat\,\iotastar\psi) + i\partial_x\psi}_{L^2} < \varepsilon.
      \]
    \end{theorem}

    \begin{proof}
      [Proof sketch]
Let $\psi_n \in \Schwartz$ with $\psi_n \to \psi$ in $\ltwo$.
By Lemma~\ref{lem:uniform-bound}, the operators $\iotaq \Xhat \iotastar$ and
$\iotaq \Phat \iotastar$ are bounded on the range of $\iotastar$ with norm $\leq q/\pi$.
A standard $3\varepsilon$ argument using
(i) the uniform bound of Lemma~\ref{lem:uniform-bound} (although the norm $\leq q/\pi$
is not uniform in $q$, this is sufficient for the density argument since $\psi_n$ is fixed
before taking $q \to \infty$),
(ii) the Schwartz convergence of Theorem~\ref{thm:schwartz} for $\psi_n$, and
(iii) the density of $\Schwartz$ in $\ltwo$,
yields the result.
    \end{proof}

  \section{Conclusion}
    We have established the strong convergence of the rescaled Weil generators
    $\Xhat$ and $\Phat$ on $\Heis(\Zq)$ to the position and momentum operators
    on $\ltwo$, thereby proving the statement formulated as hypothesis~[H2] in version~2.0
    of~\cite{Beau2026q5a}.
    The proof isolates two structural ingredients: a discrete Sobolev identity,
    which controls the local finite-difference error at the $\ell^2$ level, and a
    quantitative aliasing estimate, which governs the passage from $\ell^2$ to
    $L^2$ with an explicit rate.
    Their combination yields, in particular, a quantitative quasi-isometry for the
    sinc embedding with rate $O(q^{-1})$, whose possible use in an identification of the
    effective metric constant $A$ is discussed in Remark~\ref{rem:reusability}.
    In the framework of $Q5a$ version~2.0 this is the generator part of the large-$q$ limit.
    It supplies no spatial limit operator: for uniformly bounded weights, the current $Q5a$ proves
    that the published admissibility forms have zero-form Mosco limit (\cite{Beau2026q5a},
    Theorem~5.1;
    Remark~\ref{rem:Q5a-consequence}).

    \newpage
    \bibliographystyle{plain}
    \bibliography{cosmochrony-bibliography, external-refs}

\end{document}
